% !PRINTED 365
\weberpage{365}
%<<CONT
definirt sind, so folgt
\[
T_n S_{n-1} - S_n T_{n-1} = (-1)^n (T_0 S_1 - T_1 S_0),
\]
also im Besonderen

\weq{11}{P_n Q_{n-1} - Q_n P_{n-1} = (-1)^n.}

Diese Formel, von der wir noch mannigfache Anwendung machen werden, zeigt, dass die Zahlen $P_n$, $Q_n$ ohne gemeinsamen Theiler sind, dass also die Näherungsbrüche $P_n : Q_n$ nicht durch Heben reducirt werden können.

Ist der Kettenbruch ein endlicher, ist also die Schlusszahl $x_{n-1} = a_{n-1}$ eine ganze Zahl, so setzen wir die Bildung der Näherungsbrüche nicht weiter fort als bis zu
\[
\frac{P_n}{Q_n} = \frac{P_{n-1} a_{n-1} + P_{n-2}}{Q_{n-1} a_{n-1} + Q_{n-2}},
\]
und die Formel (2) zeigt, dass dann also der \gesperrt{letzte Näherungsbruch mit dem Werthe des Kettenbruches übereinstimmt.}

\begin{center}
§. 118.

Lösung unbestimmter Gleichungen mit zwei Unbekannten.
\end{center}

Die zuletzt gefundenen Formeln führen zur Lösung einer Aufgabe, die in vielen Anwendungen vorkommt:

\gesperrt{Es seien $\alpha$, $\beta$ zwei gegebene ganze Zahlen ohne gemeinsamen Theiler, es sollen zwei andere ganze Zahlen gefunden werden, die der Bedingung genügen:}

\weq{1}{\alpha y - \beta x = 1.}

Diese Aufgabe hat, wenn sie überhaupt lösbar ist, unendlich viele Lösungen, die alle aus einer von ihnen abgeleitet werden können.

Ist nämlich $x_0$, $y_0$ eine Lösung, also

\weq{2}{\alpha y_0 - \beta x_0 = 1,}
so folgt durch Subtraction von (1) und (2)
\[
\alpha (y - y_0) = \beta (x - x_0),
\]
woraus zu schliessen ist, da $\alpha$ mit $\beta$ keinen gemeinsamen Theiler
%CONT>>
